\centerline{\bf \Huge Построение конструктивов II}


\begin{flushright}
        \scriptsize{Единственный счастливый предмет, который вам понадобится, уже находится внутри — УВЕРЕННОСТЬ\\ Мидорима Синтаро}
\end{flushright}

\problem{\textbf{1}} Придумайте трехзначное число, запись которого состоит из различных цифр, следующих в порядке возрастания, а в названии этого трехзначного числа все три слова начинаются с одной и той же буквы.\\

\problem{\textbf{2}} Закрасьте некоторые клетки квадрата $4\times4$ так, чтобы любая закрашенная клетка имела общую сторону ровно с тремя незакрашенными.\\

\problem{\textbf{3}} Сложите квадрат $8\times8$ из 32 прямоугольников $1\times2$  так, чтобы ни в какой точке не сходились четыре прямоугольника.\\

\problem{\textbf{4}} В клетках таблицы $4\times4$ расставьте положительные и отрицательные числа так, чтобы сумма чисел, стоящих в углах любого квадрата $2\times2$, $3\times3$ или $4\times4$, была равна 0.\\

\problem{\textbf{5}} В ряд выложены карточки, на которых написаны цифры 7, 8, 9, 4, 5, 6, 1, 2, 3. Разрешается выбрать несколько карточек, лежащих подряд, и переложить их в обратном порядке. Можно ли за три таких операции добиться порядка 1, 2, 3, 4, 5, 6, 7, 8, 9? (Конечно, можно - это же конструктивы).\\


\problem{\textbf{6}} Как расположить на футбольном поле 6 футболистов, чтобы каждый из них имел возможность сделать прямолинейную передачу по земле ровно четырем другим? (передача выполняется по земле, то есть нельзя перекидывать мяч через игрока)\\

\problem{\textbf{7}} Каждую грань куба разбили на четыре одинаковых квадрата. Можно ли каждый из получившихся квадратов покрасить в один из трёх цветов так, чтобы любые два квадрата, имеющие общую сторону, были покрашены в разные цвета? (Можно или нет? Ответ очевиден)